% Figure: battery model. (a) open-circuit voltage versus state of charge,
% with a flat plateau; (b) terminal voltage under load, showing IR drop
% and the model fit against measured discharge points.
\begin{figure}[htbp]
\centering
\begin{tikzpicture}[font=\small]
% ---- panel (a): OCV vs SOC ----
\begin{scope}
\draw[->, thick, draw=engblack!70] (0,0) -- (5.4,0)
  node[right, font=\scriptsize] {state of charge};
\draw[->, thick, draw=engblack!70] (0,0) -- (0,4.0)
  node[above, font=\scriptsize] {open-circuit voltage};
% flat plateau OCV curve: steep rise, long plateau, steep drop at the ends
\draw[very thick, draw=engblue!80]
  (0.2,0.5) .. controls (0.9,2.9) and (1.4,3.1) .. (2.6,3.2)
  .. controls (3.8,3.3) and (4.3,3.4) .. (4.9,3.9);
% plateau annotation
\draw[dashed, draw=enggray] (1.4,3.15) -- (4.0,3.15);
\node[font=\scriptsize, color=enggray, anchor=south] at (2.7,3.15) {flat plateau};
% data points
\foreach \px/\py in {0.4/1.6, 0.9/2.95, 1.8/3.15, 2.8/3.22, 3.8/3.32, 4.6/3.7} {
  \filldraw[engblack!75] (\px,\py) circle (0.05);
}
\node[font=\scriptsize, anchor=north] at (2.6,-0.5) {(a) voltage plateau};
\node[font=\scriptsize, color=engred!80, anchor=west, text width=3.0cm]
  at (0.7,1.0) {steep near the ends};
\end{scope}
% ---- panel (b): terminal voltage under load, fit vs data ----
\begin{scope}[xshift=7.0cm]
\draw[->, thick, draw=engblack!70] (0,0) -- (5.4,0)
  node[right, font=\scriptsize] {discharge time};
\draw[->, thick, draw=engblack!70] (0,0) -- (0,4.0)
  node[above, font=\scriptsize] {terminal voltage};
% cutoff line
\draw[dashed, draw=engred!70] (0,0.7) -- (5.0,0.7);
\node[font=\scriptsize, color=engred!70, anchor=east] at (5.0,0.85) {cutoff};
% model fit curve: gentle sag then a knee down to cutoff
\draw[very thick, draw=engblue!80]
  (0.1,3.5) .. controls (1.8,3.0) and (3.2,2.7) .. (4.0,2.2)
  .. controls (4.4,1.7) and (4.6,1.0) .. (4.75,0.7);
% measured discharge points scattered on the fit
\foreach \px/\py in {0.3/3.45, 1.2/3.15, 2.2/2.85, 3.2/2.68, 3.9/2.28, 4.5/1.15} {
  \pgfmathsetmacro\jj{0.05*sin(\px*400)}
  \filldraw[engblack!80] (\px,{\py+\jj}) circle (0.055);
}
% range prediction marker
\draw[dashed, draw=engorange!85] (4.75,0) -- (4.75,0.7);
\node[font=\scriptsize, color=engorange!85, anchor=south west] at (4.0,0.15)
  {predicted range};
\node[font=\scriptsize, anchor=north] at (2.6,-0.5) {(b) discharge fit};
\end{scope}
\end{tikzpicture}
\caption[A battery model for range estimation]{A lumped battery model for
predicting electric-vehicle range. Panel (a): open-circuit voltage against
state of charge, with a long flat plateau that makes voltage a poor direct
gauge of remaining charge in the middle of the range and a sharp one only
near the extremes, the structural fact that forces the estimator to
integrate current (coulomb counting) rather than read voltage. Panel (b):
terminal voltage under a driving load, the sum of the open-circuit term
and an internal-resistance drop; the fitted model tracks the measured
discharge points, and the time at which the terminal voltage reaches the
cutoff sets the predicted range. The model is identified from a
controlled discharge and validated against a held-out drive cycle, the
chapter's cycle applied to a quantity a driver reads off the dashboard.}
\label{fig:vol02ch18:battery-soc}
\end{figure}
